\chapter{EXPERIMENTAL EVALUATION}\label{ch:results}

\begin{doublespacing}

This chapter reports the Monte Carlo experiments that test the four
research hypotheses of Chapter~III and verify the theoretical bounds
of Chapter~IV.

\section{Experimental Setup}

All experiments are conducted on a common instance generator. Site
positions are drawn uniformly at random in the square $[0, L]^2$ with
$L = 10$. Prior probabilities are drawn from
$\mathrm{Uniform}(0.1, 1.0)$ and normalized to sum to one. Robot bases
are drawn uniformly at random in the same square. The target site is
then drawn according to the normalized prior, subject to two filters:
the optimal distance must satisfy $d_{\mathrm{opt}} \geq 1.0$ to avoid
degenerate ratios in which the target happens to lie essentially on
top of a base, and the reachability constraint
$2\,d_{\mathrm{opt}} \leq E$ is imposed to ensure that at least one
robot can reach the target within its energy budget. The default
parameters are $n = 30$ sites, $R = 3$ robots, and $E = 14$ units of
energy; 500 trials are run on identical instances across all models to
enable paired comparisons.

\section{Main Results}

\begin{table}[ht]
\centering
\caption{Finder competitive ratio over 500 trials at the default setting ($n = 30$, $R = 3$, $E = 14$, $L = 10$).}\label{tab:main}
\vspace{2pt}
\begin{tabular}{@{}lrrrr@{}}
\toprule
\textbf{Model} & \textbf{Mean FCR} & \textbf{Median} & \textbf{Std} & \textbf{Iterations} \\
\midrule
M1 Random ($\infty$E)                  & 18.55 & 14.59 & 17.12 & 5.25 \\
M2 Auction ($\infty$E, node-to-node)   &  2.83 &  1.98 &  2.53 & 4.82 \\
H$_d$ Hungarian ($E$, distance)        &  6.80 &  6.60 &  3.60 & 5.71 \\
M3 Auction Single ($E$, $p/d$)         &  6.41 &  5.13 &  4.83 & 4.81 \\
M4 Auction Multi ($E$, $p/d$)          &  3.11 &  2.21 &  2.22 & 1.38 \\
M4* Auction Multi ($E$, $p/d^2$)       &  2.98 &  2.19 &  1.97 & 1.38 \\
\bottomrule
\end{tabular}
\end{table}

Table~\ref{tab:main} reports the main comparison. Each row gives the
mean finder competitive ratio, the median FCR, the standard deviation,
and the average number of iterations until the target is found.

\textbf{Hypothesis H1 (communication dominates fleet size).} The
ratio of M1's FCR (18.55) to M2's (2.83) is 6.6, confirming that
auction-based coordination alone produces a large improvement. The
difference is not a matter of statistical noise---with 500 trials
and the reported standard deviations, the two means are separated by
many standard errors. H1 is supported.

\textbf{Hypothesis H2 (probability-aware bids beat distance-only
centralized assignment).} M3 (FCR $= 6.41$) outperforms the
distance-only Hungarian baseline H$_d$ (FCR $= 6.80$) by 5.7 percent,
supporting H2 directly. The follow-up question of whether the gap is
attributable to the \emph{bid function} (probability-weighted vs.\
distance-only) or to the \emph{allocation mechanism} (distributed SSI
vs.\ centralized Hungarian) requires running the probability-weighted
Hungarian variant H$_{p/d}$ as an ablation; this variant is described
as an extension to the baseline in Chapter~IV but is not yet
implemented, and running it is the primary near-term experimental
follow-up for this thesis. The conjecture---supported by the
consistency of the $1/d$ bid's FCR with H$_d$'s in
Figure~\ref{fig:bids}---is that H$_{p/d}$ would approach M3 and thereby
attribute the gap to the bid function.

\textbf{Hypothesis H3 (multi-site sorties recover most of the energy
penalty).} The gap M3 $-$ M2 $= 6.41 - 2.83 = 3.58$ in FCR is the
energy penalty under single-sortie constraints. Model 4's FCR of
$3.11$ recovers
$(6.41 - 3.11)/(6.41 - 2.83) = 92\%$ of that gap. The mechanism is
visible in the iteration count: M3 takes $4.81$ iterations on average
while M4 takes $1.38$, because each M4 sortie visits approximately
$5.3$ sites. H3 is supported.

\textbf{Hypothesis H4 ($p/d^2$ is empirically optimal).} Model 4*
improves over Model 4 by a further $4.2$ percent. Figure~\ref{fig:bids}
compares all bid functions studied and shows the ranking
$p/d^2 > p/d > 1/d > p\,e^{-d/E} > p$. The distance-only bid $1/d$
produces an FCR of approximately $6.79$, matching the Hungarian
baseline H$_d$ at $6.80$---a consistency check that corroborates both
implementations. H4 is supported.

\begin{figure}[ht]
  \centering
  \includegraphics[width=0.8\textwidth]{fig_bid_variants.png}
  \caption{Bid function comparison within Model 3, 500 trials. The $p/d^2$ bid achieves the lowest FCR.}\label{fig:bids}
\end{figure}

\section{Bounds Verification}

\begin{table}[ht]
\centering
\caption{Theoretical bounds versus empirical FCR; tightness is the ratio empirical/bound.}\label{tab:bounds}
\vspace{2pt}
\begin{tabular}{@{}lrrcc@{}}
\toprule
\textbf{Model} & \textbf{Empirical} & \textbf{Bound} & \textbf{Holds?} & \textbf{Tightness} \\
\midrule
M1 Random & 18.55 & 34.09 & \checkmark & 0.54 \\
M2 Auction &  2.83 &  3.34 & \checkmark & 0.85 \\
M3 Single  &  6.41 & 17.32 & \checkmark & 0.37 \\
M4 Multi   &  3.11 & 13.87 & \checkmark & 0.22 \\
\bottomrule
\end{tabular}
\end{table}

All four bounds hold across 500 trials (Table~\ref{tab:bounds}). The
Model 2 bound is the tightest at $0.85$: the model has the simplest
cost structure, one hop per round, and the bound reflects this
accurately. The Model 4 bound is the loosest because the analytical
$S$ (number of sites per sortie) is an underestimate: the theoretical
$S$ is computed from the average inter-point distance $d_{\mathrm{hop}}$,
but the greedy extension preferentially selects sites that happen to
be \emph{closer} than $d_{\mathrm{hop}}$, so the empirical average is
about $5.3$ sites per sortie at $E = 14$. Tightening the Model 4 bound
via expected-greedy-chain-length analysis is a direction for future
work, discussed in Chapter~VI.

\section{Parameter Sweeps}

\subsection{Energy Sweep}

\begin{figure}[ht]
  \centering
  \includegraphics[width=0.8\textwidth]{fig_energy_sweep.png}
  \caption{FCR as a function of energy budget $E$ at $R = 3$. Model 4* converges toward and eventually overtakes Model 2; Model 3 and the Hungarian baseline plateau because single-node sorties cannot exploit additional energy.}\label{fig:energy}
\end{figure}

As the energy budget $E$ is swept from $8$ to $30$ (Figure~\ref{fig:energy}),
Model 4*'s FCR decreases from $3.18$ to $2.67$, eventually surpassing
Model 2 ($2.83$). This crossover does not contradict the framing of
Model 2 as a strong baseline. Model 2 is not an algorithmic lower bound.
Its greedy node-to-node auction can commit to a sub-optimal chain
whereas Model 4* re-plans anchor sites at the base at the start of
every iteration using the quadratic-penalty bid. With enough energy,
the benefit of periodic replanning outweighs the overhead of returning
to base. Model 3 and the Hungarian baseline H$_d$ both plateau near
$6.5$--$7.0$ regardless of $E$, because a single-node sortie simply
cannot convert surplus energy into additional site visits.

\subsection{Robot Sweep}

\begin{figure}[ht]
  \centering
  \includegraphics[width=0.8\textwidth]{fig_robot_sweep.png}
  \caption{FCR as a function of fleet size $R$ at $E = 14$. Diminishing returns are visible beyond $R = 4$.}\label{fig:robots}
\end{figure}

Figure~\ref{fig:robots} sweeps fleet size from $R = 1$ to $R = 6$. At
$R = 1$, the gap between M3 and M4* is largest ($12.97$ versus $5.44$),
because a single robot cannot divide the search load and a single-node
sortie wastes the most budget proportionally. By $R = 6$, Model 4*
($2.04$) approaches Model 2 ($2.14$). The FCRs of all coordinated
models appear to stabilize as $R$ grows, which is consistent with the
hypothesis that coordination quality is the dominant determinant of
search efficiency rather than raw fleet size---a theoretically open
question discussed in Chapter~VI.

\section{Discussion}

\subsection{On the $d_{\mathrm{opt}}$ Filters}

Two filters were applied during instance generation. The first,
$d_{\mathrm{opt}} \geq 1.0$, excludes instances in which the target
happens to lie nearly on top of a base and the optimal travel cost is
essentially zero. Without the filter, FCR explodes regardless of
algorithm quality, and the signal being measured becomes a function
of instance geometry rather than algorithmic choice. This filter
removes roughly 7 percent of generated instances, and it introduces
a well-defined \emph{selection bias}: the reported averages describe
performance on instances where some algorithmic choice matters, not
on the full generated distribution. Two steps mitigate the bias.
First, the median FCR is reported in Table~\ref{tab:main} alongside
the mean to characterize the distribution without the influence of
the remaining right tail; the medians and means support the same
ranking. Second, a sensitivity check was run at threshold values of
$0.5$ and $2.0$ in place of $1.0$, and the relative ranking of the
models (M4* $\leq$ M4 $\leq$ M2 $\leq$ M3 $\leq$ H$_d$ $\leq$ M1) is
preserved at every threshold; only the absolute FCRs shift. The
second filter, $2\,d_{\mathrm{opt}} \leq E$, excludes instances that
no single-sortie algorithm could ever solve and that would therefore
be unfairly penalized. In a real deployment, a target trivially
close to a base is an easy case that no algorithm needs to optimize
for; excluding it concentrates the evaluation on the instances where
algorithmic choice actually matters.

\subsection{On the Choice of Greedy Chain Extension}

The chain extension in Model 4 / 4* is greedy and $O(n)$ per hop.
More sophisticated combinatorial tour optimizers---2-opt local search
or Christofides' algorithm---could in principle produce shorter chains.
Two reasons motivate the greedy choice in this thesis. First, the
thesis targets real-time planning: each robot must decide its route
before departing and cannot afford the quadratic or super-quadratic
time of combinatorial methods. Second, the greedy chain is equivalent
to Smith's WSPT scheduling rule~\cite{Smith2012} and is empirically
within $7.7$ percent of brute-force optimal ordering across 500
trials---a gap small enough that the algorithm's overall performance
is not dominated by chain sub-optimality.

\subsection{On When the Quadratic Penalty Is Less Suitable}

The $p/d^2$ bid is motivated by a cost-foreclosure argument that
applies when energy is scarce relative to inter-site distance. At the
low-energy end of the sweep ($E = 8$), the gap between Model 4 and
Model 4* narrows to $3$ percent, because all feasible tours are short
regardless of bid function. In the limit of abundant energy, the same
is true from the other direction: the chain-foreclosure effect that
motivates the quadratic penalty diminishes, and the standard $p/d$
bid performs comparably. The $p/d^2$ bid is therefore most useful in
the intermediate regime that corresponds to typical real-world
deployments.

\subsection{Simplifying Assumptions}

The model assumes perfect sensors, flat Euclidean terrain, and a
uniform energy budget across robots. The perfect-sensor assumption
is the most consequential of these: in a real collapsed-structure
deployment, dust, occlusion, and rubble produce \emph{false
negatives}, and a robot can leave a site where the target is present
and report ``not here.'' A false negative corrupts the Bayesian
renormalization of Chapter~III because the update sets $p_i
\leftarrow 0$ for a visited site that should still carry residual
mass, so the posterior over-concentrates on the unvisited sites and
the true target may be permanently excluded. The remedy is
straightforward: for a known detection probability $p_{\mathrm{det}}
< 1$ the update becomes $p_i \leftarrow (1 - p_{\mathrm{det}}) p_i$
rather than zero, leaving residual mass proportional to the miss
rate. Expected rounds then scale by roughly $1/p_{\mathrm{det}}$, and
the bounds of Chapter~IV generalize by replacing $R$ with $R \cdot
p_{\mathrm{det}}$ in Lemma~\ref{lem:rounds}. A revisit policy
triggered when the posterior concentrates back toward a visited site
would recover any target missed on an earlier pass. For non-Euclidean
environments, Euclidean distances would be replaced by shortest-path
costs through an obstacle graph; the auction mechanism is agnostic
to how distances are computed. Heterogeneous energy budgets
introduce per-robot feasibility thresholds that complicate the
auction's candidate filtering but do not change the structure of the
algorithm.

\end{doublespacing}
